% Originally: content/week-04.tex, \section{The differential} (Corsin Week 4, Monday)

% Source: Corsin Nick/Class Notes/Week 4.pdf, p. 2
\section{The differential}
\label{sec:the_differential}

For $f : \mathbb{R}\to\mathbb{R}$ differentiable at $x_0 \in \mathbb{R}$:
\[f(x_0+h) = f(x_0) + f'(x_0)h + o(h)\]
\[\Updownarrow\]
\[\lim_{h\to 0}\frac{f(x_0+h)-f(x_0)}{h} = f'(x_0)\]
\[\Updownarrow\]
\[\lim_{h\to 0}\left|\frac{f(x_0+h)-f(x_0)-f'(x_0)h}{h}\right| = 0\]

For $F : \mathbb{R}^n \to \mathbb{R}^m$ differentiable at $x_0 \in \mathbb{R}^n$:
\[F(x_0+h) = F(x_0) + \underbrace{DF_{x_0}(h)}_{\text{linear map}} + o(|h|)\]
\[\Updownarrow\]
\[\lim_{h\to 0}\frac{|F(x_0+h)-F(x_0)-DF_{x_0}(h)|}{|h|} = 0\]
\[\Downarrow\]
\[\lim_{s\to 0}\frac{F(x_0+sv)-F(x_0)}{s} = DF_{x_0}(v) = \partial_vF(x_0) \quad \text{for each fixed } v \in \mathbb{R}^n\]
\[\text{(the \newterm{directional derivative} of } F \text{ at } x_0 \text{ in direction } v\text{)}\]

Note that the last arrow points \textbf{one way only}. The first two lines say the same thing;
the third is a strictly weaker consequence, and recovering differentiability from it is exactly
what fails in several variables; see \cref{ex:partial_derivatives_not_continuous} below.

% Transition: Claude Opus 5 (High)
It is worth pausing on why the definition had to be restated rather than merely copied. The
one-variable derivative is defined as a \emph{quotient}, and that is precisely what does not
survive: for $h \in \mathbb{R}^n$ with $n \geq 2$ the expression
$\big(F(x_0+h)-F(x_0)\big)/h$ is meaningless, because one cannot divide by a vector. The first
line above is the version that does survive, and it is the one to remember: \textbf{$F$ is
differentiable at $x_0$ if it is approximated near $x_0$ by an affine map, to an error smaller
than $|h|$.} Everything that follows, the Jacobian, the chain rule and Taylor expansions alike,
is bookkeeping for that single idea. The reason the derivative becomes a \emph{linear map} rather
than a number is simply that a linear map is what an approximation in $n$ variables looks like.

% Generator: Claude Opus 5.5 (Extra) --- what "error o(|h|)" looks like, in one variable.
% f(x) = 1 + 0.5x + 0.35x^2 near x_0 = 0, tangent 1 + 0.5x. For eps = 0.35 the wedge between
% the lines of slope 0.5 - 0.35 = 0.15 and 0.5 + 0.35 = 0.85 through (0,1) contains the graph
% exactly when 0.35x^2 <= 0.35|x|, i.e. |x| <= 1; so delta = 1, and the graph leaves the wedge at
% (1, 1.85) and (-1, 0.85), both of which lie on a wedge edge (1 + 0.85 = 1.85, 1 - 0.15 = 0.85).
% Wedge polygon over [-2.2, 2.2]: right half between heights 1.33 and 2.87 at x = 2.2, left half
% between -0.87 and 0.67 at x = -2.2. A smaller eps gives a thinner wedge and a smaller delta.
\begin{center}
\begin{tikzpicture}[>=Latex, scale=1.0,
  lbl/.style={fill=white, fill opacity=0.85, text opacity=1, inner sep=1.2pt, rounded corners=1pt}]
  \fill[orange!18] (0,1) -- (2.2,1.33) -- (2.2,2.87) -- cycle;
  \fill[orange!18] (0,1) -- (-2.2,0.67) -- (-2.2,-0.87) -- cycle;
  \draw[orange!80!black, dashed] (-2.2,-0.87) -- (2.2,2.87);
  \draw[orange!80!black, dashed] (-2.2,0.67) -- (2.2,1.33);
  \draw[->, gray] (-2.5,-1.2) -- (2.6,-1.2) node[right, font=\scriptsize, text=black] {$x$};
  \draw[thick, OliveGreen] (-2.3,-0.15) -- (2.3,2.15);
  \draw[very thick, blue!70!black, domain=-2.2:2.2, samples=60, smooth]
    plot (\x, {1 + 0.5*\x + 0.35*\x*\x});
  \fill (0,1) circle (1.6pt);
  \fill[blue!70!black] (1,1.85) circle (1.3pt);
  \fill[blue!70!black] (-1,0.85) circle (1.3pt);
  \draw[gray, dotted] (1,-1.2) -- (1,1.85);
  \draw[gray, dotted] (-1,-1.2) -- (-1,0.85);
  \draw[gray] (0,-1.15) -- (0,-1.25) node[below, font=\scriptsize, text=black] {$x_0$};
  \draw[<->, thick] (-1,-0.95) -- (1,-0.95);
  \node[lbl, font=\scriptsize] at (0.5,-0.95) {$\lvert h\rvert<\delta$};
  \node[lbl, font=\scriptsize, text=blue!70!black] at (-1.65,2.0) {graph of $f$};
  \node[lbl, font=\scriptsize, text=OliveGreen, anchor=west] at (2.35,2.15) {tangent: slope $f'(x_0)$};
  \node[lbl, font=\scriptsize, text=orange!70!black, anchor=west, align=left] at (2.35,1.4)
    {wedge: slopes\\ $f'(x_0)\pm\varepsilon$};
\end{tikzpicture}
\end{center}
For every $\varepsilon > 0$ there is a $\delta > 0$ such that, over $\lvert h\rvert < \delta$, the
graph lies inside the shaded wedge around the tangent line. That is the statement
$\lvert f(x_0+h)-f(x_0)-f'(x_0)h\rvert \leq \varepsilon\lvert h\rvert$, and it is the form that
generalises: in several variables the line becomes the graph of the affine map
$h \mapsto F(x_0) + DF_{x_0}(h)$, and the wedge becomes a cone around it.

\begin{ainote}
Corsin's middle line reads
$\lim_{h\to0}\frac{F(x_0+h)-F(x_0)}{|h|} = DF_{x_0}\!\left(\frac{h}{|h|}\right)$, inside a chain
of $\Updownarrow$. As written that limit does not exist: $h$ is being sent to $0$ on the left
while still appearing on the right, and the quotient genuinely depends on the direction of
approach. The intended statement is the one typeset above, in which a direction $v$ is fixed first
and only then the scalar $s$ is sent to $0$, and it is only an \emph{implication}, not an
equivalence. The distinction is not pedantic: it is the entire content of
\Cref{ex:partial_derivatives_not_continuous} and of
\Cref{rem:directional_derivative_caveat}.
\end{ainote}

% Source: Corsin Nick/Class Notes/Week 4.pdf, p. 2
% (Re-cited: the paragraph and AI-Note above are additions, not transcription.)
\begin{definition}[Differentiability]
\label{def:differentiable}
We say that $F : \mathbb{R}^n \to \mathbb{R}^m$ is \newterm{differentiable}
\germanfor{differenzierbar} at $x_0$ if a linear map $DF_{x_0}$ as above exists, where
$x_0, h \in \mathbb{R}^n$. The same definition applies verbatim to $F : U \to \mathbb{R}^m$
for any open subset $U \subseteq \mathbb{R}^n$, and we use it in that generality throughout.
\end{definition}

% Sonnet 5 (Medium)
\begin{remark}
The linear map $DF_{x_0}$ in \cref{def:differentiable}, when it exists, is automatically unique:
if $L_1, L_2$ both satisfy the defining limit, then
$(L_1-L_2)(h)/|h| \to 0$ as $h \to 0$ along every direction, which forces the linear map
$L_1 - L_2$ to vanish identically. Concretely: write $L := L_1-L_2$ and suppose $L(h) \neq 0$
for some $h$. Along the ray $th$, $t > 0$, linearity gives
\[\frac{|L(th)|}{|th|} = \frac{t\,|L(h)|}{t\,|h|} = \frac{|L(h)|}{|h|} \neq 0,\]
a nonzero \emph{constant}. Consequently, the quotient cannot tend to $0$ as $t \to 0$, and
therefore $L(h)=0$ for every $h$. This
is the multivariable analogue of the uniqueness-of-limits argument from
\Cref{prop:uniqueness_of_limits}.
\end{remark}

% Supplement: lec_notes.pdf, p. 37 (Lemma 10.9)
% Generator: Claude Opus 5 (High) --- used silently throughout the document (every time a
% Jacobian's rows are read off as the gradients of the components, and in the proof of the
% sufficient condition for differentiability, which reduces to m = 1) but never stated.
\begin{lemma}[Differentiability is componentwise]
\label{lem:differentiability_componentwise}
Let $U \subseteq \mathbb{R}^n$ be open, $F : U \to \mathbb{R}^m$, and write
$F = (F_1,\dots,F_m)$ for its components $F_j : U \to \mathbb{R}$. Then $F$ is differentiable at
$x \in U$ if and only if every $F_j$ is, and in that case
\[\big(DF_x(v)\big)_j = D(F_j)_x(v) \qquad \text{for all } v \in \mathbb{R}^n \text{ and } 1 \leq j \leq m.\]
\end{lemma}

\begin{proof}
Both directions come from the same observation: a vector in $\mathbb{R}^m$ is $o(|h|)$ exactly
when each of its $m$ components is, since
$\max_j|w_j| \leq |w| \leq \sqrt{m}\max_j|w_j|$.

If each $F_j$ is differentiable at $x$ with differential $L_j$, assemble
$L(v) := (L_1(v),\dots,L_m(v))$, which is linear. The remainder of $F$ at $x$ has $j$-th component
the remainder of $F_j$, each of which is $o(|h|)$, hence so is the vector. Conversely, if $F$ is
differentiable with differential $L$, then composing with the $j$-th coordinate projection, which
is linear and $1$-Lipschitz, shows $F_j$ is differentiable with differential $v \mapsto L(v)_j$.
\end{proof}

\begin{remark}[What this buys, and what it does not]
\label{rem:componentwise_scope}
This is the licence to reduce to $m = 1$, and it is used in exactly that way twice below: the
rows of the Jacobi matrix are the gradients of the components (\cref{sec:jacobi_matrix}), and the
proof that continuous partial derivatives suffice
(\cref{thm:continuous_partials_imply_differentiable}) treats a scalar function only.

It does \emph{not} say anything about the other index. Differentiability in $n$ variables does
not reduce to differentiability in each variable separately: that is precisely the failure
\cref{ex:partial_derivatives_not_continuous} exhibits. The target $\mathbb{R}^m$ splits; the
source $\mathbb{R}^n$ does not.
\end{remark}

\begin{proposition}[Differentiability implies continuity]
\label{prop:differentiable_implies_continuous}
If $F : U \to \mathbb{R}^m$, $U \subseteq \mathbb{R}^n$ open, is differentiable
(\cref{def:differentiable}) at $x_0 \in U$, then $F$ is continuous at $x_0$.
\end{proposition}

\begin{proof}
By \cref{def:differentiable}, $F(x_0+h) - F(x_0) = DF_{x_0}(h) + o(|h|)$ as $h \to 0$. Since
$DF_{x_0}$ is linear on the finite-di\-men\-sional space $\mathbb{R}^n$, it is bounded, i.e.\ there
exists $C \geq 0$ with $|DF_{x_0}(h)| \leq C|h|$ for all $h$. Hence
\[|F(x_0+h) - F(x_0)| \leq |DF_{x_0}(h)| + o(|h|) \leq C|h| + o(|h|) \xrightarrow{h\to 0} 0,\]
so $F(x_0+h) \to F(x_0)$ as $h \to 0$, i.e.\ $F$ is continuous at $x_0$.
\end{proof}

\begin{remark}
The converse fails: continuity does not imply differentiability, exactly as in one variable
(e.g.\ $x \mapsto |x|$ at $0$). \Cref{ex:partial_derivatives_not_continuous} below shows the
gap is even wider in several variables, where \emph{all} partial derivatives can exist at a point
without $F$ even being continuous there, let alone differentiable.
\end{remark}
